\cotefile{64}
\scanpage{145}
classes d'isomorphie d'objets \unsure{d'icelles|d'icelle} est en corr. biunivoque canonique avec l'ens des applications holomorphes
\makeatletter\tagsleft@true\makeatother
\[
\tau:S\longrightarrow\mathbb C\setminus\mathbb R . \tag*{(5)}
\]
On dit qu'une structure de Torelli est ``directe'' ou ``stricte'' si \deleted{\unsure{l'est|l'ori.}} pour $\unsure{\forall|\mathrm{tt}}\,s\in S$, la base \deleted{de} $z_1,z_2$ de $\unsure{H|X|\mathcal H}_s\otimes_{\mathbb Z}\mathbb R$ \deleted{\unsure{correspond à|correspond|correspond.}}\inserted{définit} l'\deleted{\unsure{orient.|struc.}}\inserted{orientation} de cet espace correspondant à sa structure complexe, i.e si (5) est un morphisme
\[
\tau:S\longrightarrow\mathcal D^+=\{z\in\mathbb C\mid\unsure{J|\mathrm I|\Im}z>0\}. \tag*{(6)}
\]
Donc $\mathcal D^+$ apparaît comme l'espace analytique modulaire pour la structure ``courbe elliptique relative avec str. de Torelli stricte sur espace analytique variable'', la famille modulaire étant donnée par
\[
X_{\mathcal D^+}=(\mathcal D^+\times\mathbb C)/\mathbb Z^2 \tag*{(7)}
\]
où $\mathbb Z^2$ opère sur $\mathcal D^+\times\mathbb C$ par
\[
\begin{aligned}
\deleted{\unsure{m,n}}\;T_{m,n}(\tau,z)&=(\tau,z+m+n\tau)\\
&\hspace{5em}\tau\in\mathcal D^+,\ z\in\mathbb C
\end{aligned}\tag*{(8)}
\]
de sorte que l'on a
\[
\begin{aligned}
(X_{D^+})_\tau &\simeq \mathbb C/(\mathbb Z+\mathbb Z.\tau)\simeq\mathbb C^*/\text{sous-groupe}\\
&\hspace{10em}\text{engendré}\\
&\hspace{10em}\text{par }\deleted{\unsure{q}}\;k=\exp(2i\pi\tau)\\
&\hspace{10em}\text{(\unsure{satisfaisant|on trouve bien} }|k|<1\text{)}
\end{aligned}\tag*{(8)}
\]
% archive sheet number: 144
\scanpage{146}
Sur l'ens des \deleted{\unsure{structu|structures}} classes d'isomorphie de structures de Torelli directes, le groupe
\[
\deleted{\unsure{SL(2,}}
\]
\[
SL(2,\mathbb Z)=\left\{\begin{pmatrix}a&b\\c&d\end{pmatrix}\;\middle|\;
\begin{aligned}a,b,c,d&\in\mathbb Z\\ad-bc&=1\end{aligned}\right\} \tag*{(9)}
\]
opère à droite, par
\[
t\unsure{\mathrel{\vdash\!\dashv}|\longmapsto}t\circ g,
\]
on va expliciter cette opération, en tant qu'opération sur $\mathcal D^+$, et $\unsure{\overline{\mathrm m}|=}$ sur $X_{\mathcal D^+}$ (vu que l'on a une opération de $SL(\unsure{2|1},\mathbb Z)$ sur la situation modulaire).

On voit qu'à une structure de Torelli $(E,(z_1,z_2))$, $g$ associe la structure $(E,(z'_1,z'_2))$, avec
\[
\begin{aligned}
z'_1&=az_1+cz_2\\
z'_2&=bz_1+dz_2
\end{aligned}
\]
d'où, si $\tau,\tau'$ sont les invariants correspondants
\[
-\qquad \tau=z_2/z_1,\quad\tau'=z'_2/z'_1
\]
la formule
\[
\tau'=\frac{d\tau+b}{\deleted{\unsure{b}}c\tau+a}\qquad
\left(\tau'=\tau\circ g,\ g=\begin{pmatrix}a&b\\c&d\end{pmatrix}\right) \tag*{(10)}
\]
Si on remplace l'opération envisagée par l'opération à gauche correspondante, \deleted{\unsure{on|---}} $t\longmapsto t\circ g^{-1}$, notée aussi $\tau\longmapsto\tau\circ g^{-1}$, on trouve,
% archive sheet number: 145
\scanpage{147}
Au niveau de la famille modulaire, l'opération qui correspond à $t\longmapsto t\circ g$ se calcule aisément, si on a
\[
\begin{aligned}
E_\tau &=\mathbb C/(\mathbb Z+\mathbb Z\tau)\\
E_{\tau\circ g}&=\mathbb C/(\mathbb Z+\mathbb Z\tau')\qquad \tau'=\tau\circ g=\frac{d\tau+b}{c\tau+a}
\end{aligned}
\]
on note que l'on a une application \inserted{$\mathbb C$} linéaire naturelle $\unsure{\varphi|\phi}\quad\mathbb C\unsure{-|\to}\mathbb C$, qui transforme $\mathbb Z+\mathbb Z\tau$ en $\mathbb Z+\mathbb Z\tau'$, qui est la multiplication par $\unsure{\zeta|\xi}=z_1/z'_1=\dfrac{z_1}{az_1+cz_2}=\dfrac{1}{a+c\tau}$,
\diagram{\begin{tikzcd}[ampersand replacement=\&,column sep=large,row sep=large]
\mathbb C \arrow[r,"\sim" {yshift=-1pt},"{\lambda\mapsto\lambda z_1,\ \substack{1\mapsto z_1\\\tau\mapsto z_2}}" {yshift=7pt}] \arrow[d,"{\lambda\mapsto\lambda\frac{z_1}{z'_1}}"',"{\unsure{\wr|2}}"] \& t_E \arrow[d,"{\Vert}"]\\
\mathbb C \arrow[r,"{\lambda\mapsto\lambda z'_1,\ \substack{1\mapsto z'_1\\\tau'\mapsto z'_2}}"' ] \& t_E
\end{tikzcd}}
\margin{NB on a $\left\{\begin{aligned}1&=a\zeta+c\zeta\tau\\\tau'&=b\zeta+d\zeta\tau\end{aligned}\right.$}
% slanted
et par passage aux quotients il s'en déduit
\[
E_{\tau\deleted{\unsure{\circ}}}\xrightarrow{\sim}E_{\tau\circ g}
\]
remplaçant $g=\begin{pmatrix}a&b\\c&d\end{pmatrix}$ par $\tau_\infty(g^{-1})=\begin{pmatrix}d&b\\c&a\end{pmatrix}$, on trouve
\[
E_\tau\xrightarrow[\substack{\text{induit par}\\\text{multiplication}\\\text{par }(c\tau+d)^{-1}}]{\sim}E_{\tau'\deleted{\unsure{=\tau}}}
\qquad \tau'=\tau\circ\tau_\infty(g^{-1})=\frac{a\tau+b}{c\tau+d}
\]
donc on fera opérer $g\in SL(2,\mathbb Z)$ \inserted{à gauche} sur $X_{\mathcal D^+}\simeq\deleted{\unsure{\mathbb Z}}(\mathcal D^+\times\mathbb C)/\mathbb Z\times\mathbb Z$ via l'action
\[
\begin{aligned}
\deleted{\unsure{T_g(\tau}}\;T_g.(\tau,z)&=\left(g.\tau,\deleted{\unsure{z}}\frac{z}{c\tau+d}\right)
&\text{où }g&=\begin{pmatrix}a&b\\c&d\end{pmatrix}\\
&&g.\tau&=\frac{a\tau+b}{c\tau+d}
\end{aligned}\tag*{(14)}
\]
% archive sheet number: 146
\scanpage{148}
vu que
\[
g^{-1}=\begin{pmatrix}d&-b\\-c&a\end{pmatrix}
\]
la formule
\[
\tau\circ g^{-1}=\frac{a\tau-b}{-c\tau+d}. \tag*{(11)}
\]
Introduisons l'élément
\[
\left\{\begin{aligned}
&\tau_\infty=\begin{pmatrix}-1&0\\0&+1\end{pmatrix}\in\unsure{Gl|GL}(2,\mathbb Z)\qquad\text{d'où}\\
&\tau_\infty\begin{pmatrix}A&B\\C&D\end{pmatrix}\tau_\infty^{-1}
=\tau_\infty\left(\begin{pmatrix}A&B\\C&D\end{pmatrix}\right)
=\begin{pmatrix}A&-B\\-C&D\end{pmatrix}\\
&\tau_\infty\left(\begin{pmatrix}a&b\\c&d\end{pmatrix}^{-1}\right)
=\begin{pmatrix}d&b\\c&a\end{pmatrix}\qquad\substack{\text{``échange de}\\\text{$a$ et de $d$''}}
\end{aligned}\right.,\tag*{(12)}
\]
on trouve
\[
\begin{aligned}
\tau\circ\tau_\infty(g^{-1})&=\frac{a\tau+b}{c\tau+d}
&\text{où }\ g&=\begin{pmatrix}a&b\\c&d\end{pmatrix}\\
&&\tau_\infty(g^{-1})&=\begin{pmatrix}d&b\\c&a\end{pmatrix}
\end{aligned}\tag*{(13)}
\]
Par la suite, quand on fait opérer $SL(2,\mathbb Z)$ sur $\mathcal D^+$, ce sera tjrs (sauf mention expresse) par la formule (13) -- ce n'est donc pas l'action ``modulaire'' la plus évidente, mais celle consacrée par l'usage, qui correspond à la précédente \inserted{(à droite)}, modifiée via le passage
\[
\underset{\substack{\Vert\\\left(\begin{smallmatrix}a&b\\c&d\end{smallmatrix}\right)}}g\longmapsto\tau_\infty(g^{-1})=\begin{pmatrix}d&b\\c&a\end{pmatrix}.
\]
% archive sheet number: 147
\scanpage{149}
de $SL(2,\mathbb Z)$ sur $\mathcal D^+\times\mathbb C$. On peut en fait faire opérer $\mathbb Z^2.\underbrace{SL(2,\mathbb Z).\deleted{\mathbb Z^2}}_{\text{$\frac12$-direct}}$ sur $\mathcal D^+\times\mathbb C$, par
\[
T_{g,m,n}(\tau,z)=\left(g\tau,\frac{z}{c\tau+d}+m+ng\tau\right) \tag*{(15)}
\]
et l'opération déduite de $SL(2,\mathbb Z)$ sur $X_{\mathcal D^+}$ se déduit de (15) en passant au quotient \deleted{\unsure{sur|sous}} dans $\mathcal D^+\times\mathbb C$ par le \deleted{\unsure{g|groupe|faux}}sous-groupe $\mathbb Z^2$ de $\mathbb Z^2.SL(2,\mathbb Z)$.

On peut énoncer dès maintenant

\uline{Théorème} La famille $\mathbb C$-analytique \unsure{de|des} courbes elliptiques, \deleted{\unsure{défini|définí|définie}}\inserted{modulaire} pour la rigidification de Torelli stricte, définie par
\diagram{\begin{tikzcd}[ampersand replacement=\&,row sep=large,column sep=small]
X_{D^+}=D^*\times\mathbb C/\mathbb Z^2 \arrow[d,"{(\tau,\dot z)\mapsto z}"] \& \substack{\mathbb Z^2\ \text{opérant par}\\T_{m,n}(\tau,z)=(\tau,z+m+n\tau)}\\
D^+ \&
\end{tikzcd}}
est un revêtement universel de $M_{1,1}$, avec \deleted{(i.e. \unsure{elle|il} est localement modulaire, et \unsure{épimorphique|épimorphisme},} comme groupe d'automorphismes sur $M_{1,1}$ le groupe $SL(2,\mathbb Z)$, opérant par
\[
g\tau=\frac{a\tau+b}{c\tau+d}\qquad g=\begin{pmatrix}a&b\\c&d\end{pmatrix},\qquad\begin{aligned}a,b,c,d&\in\mathbb Z\\ad-bc&=1\end{aligned}
\]
L'opération précédente est aussi décrite, en termes du pb modulaire de Torelli, par
% archive sheet number: 148
\scanpage{150}
\[
\begin{aligned}
g\cdot\tau=\tau\circ\tau_\infty(g^{-1})\qquad \text{où}\quad
&\tau_\infty=\begin{pmatrix}-1&0\\0&1\end{pmatrix}\in\mathrm{Gl}(2,\mathbb Z)\\
&\tau_\infty(g^{-1})=\begin{pmatrix}d&b\\c&a\end{pmatrix}\quad \text{si}\\
&\hspace{5em}g=\begin{pmatrix}a&b\\c&d\end{pmatrix}
\end{aligned}
\]

\uline{Corollaire 1} Le groupe fondamental $\unsure{\mathcal T|\mathscr T}_{1,1}$ de $M_{11}$ relatif au revêtement universel précédent est canoniquement isomorphe à $SL(2,\mathbb Z)$ (avec l'opération précédente \inserted{à gauche} de $SL(2,\mathbb Z)$ sur ledit revêtement, déduite de l'opération canonique \inserted{\unsure{à|d}} droite par le passage de $g$ à $g^{-1}$ .-- ).

Les marteaux-pilons sur les $\mathcal T_{g,\unsure{\nu|v}}$ donnent alors:

\uline{Corollaire 2} Soit $(\deleted{\unsure{\Sigma|Z|E}}X,o)$ une surface orientée compacte de genre 1, $\mathcal T$ le groupe des classes d'isotopie d'automorphismes de $(X,o)$, $\deleted{\unsure{d|o}}H=\pi_1(X,o)=H_1(X,\mathbb Z)$, \deleted{\unsure{considérons}}\inserted{alors} \unsure{l'|le} homomorphisme canonique
\[
\mathcal T\longrightarrow\mathrm{Aut}^{+}_{\mathbb Z}(H)\tag*{(16)}
\]
est un isomorphisme

\uline{Corollaire 3} Soient $\pi$ un groupe à bouts du type $(1,1)$, $H=\pi_{\mathrm{ab}}$, $\mathcal T$ le groupe des autom. ext à bouts de $\pi$ compatibles avec l'orientation, alors l'hom canonique
\[
\mathcal T\longrightarrow\mathrm{Aut}^{+}_{\mathbb Z}(H)\qquad\text{est un iso.}
\]
% archive sheet number: 149
\scanpage{151}
Considérons maintenant des courbes elliptiques \inserted{$E$} avec ``rigidification différentielle'', définie par une base $\unsure{\varphi|\phi}$ de $t_E$, permettant donc d'identifier $t_E$ à $\mathbb C$. La donnée d'une courbe elliptique avec rigidification différentielle s'identifie donc à la donnée d'un \uline{réseau} $H\subset\mathbb C$ dans $\mathbb C$. La donnée d'une courbe elliptique avec rigidification double Torelli stricte $+$ différentielle s'identifie donc à la donnée de
\[
z_1,z_2\in\mathbb C^*\quad\text{avec}\quad
\deleted{\unsure{z_2|z}}\ \tau=z_2/z_1\in\mathcal D^+\tag*{(17)}
\]
i.e. \deleted{la $\mathbb R$ base de $\mathbb C$ \unsure{d}}\inserted{$z_1,z_2$ forment une} $\mathbb R$-base de $\mathbb C$, qui est \uline{directe} pour l'orientation canonique de $\mathbb C$. On fait opérer $SL(2,\mathbb Z)$ à gauche sur \unsure{la|leur} variété analytique complexe
\[
PD^+=\{(z_1,z_2)\mid z_1,z_2\in\mathbb C^*,\ z_2/z_1\in D^+\}\tag*{(18)}
\]
par
\[
g\,(z_1,z_2)=(dz_1+cz_2,\ bz_1+az_2)\tag*{(\deleted{18}19)}
\]
% slanted
\margin{c'est donc, via le passage $g\mapsto\tau_\infty(g^{-1})$, l'opération naturelle qui découle du Pb modulaire}

La famille modulaire bi-rigidifiée $\unsure{\mathcal X|X|\mathfrak X}_{P\unsure{\Delta|D|\mathfrak D}^+}$ peut s'écrire
% archive sheet number: 150
\scanpage{152}
\[
\unsure{X|\mathcal X|\mathfrak X}_{P\mathcal D^+}\simeq PD^+\times\mathbb C/\mathbb Z^2\tag*{(20)}
\]
où $\mathbb Z^2$ opère sur $PD^+\times\mathbb C$ par
\[
T_{m,n}(z_1,z_2,z)=(z_1,z_2,z+m\deleted{\unsure{Z|z}}z_1+nz_2)
\]

Alors le groupe $\mathbb C^*\times SL(2,\mathbb Z)$ opère sur $X_{PD^+}\longrightarrow PD^+$, en faisant opérer $SL(2,\mathbb Z)$ par
\[
T_g(z_1,z_2,z)=(dz_1+cz_2,bz_1+az_2,z)\tag*{(21)}
\]
et $\mathbb C^*$ \deleted{\unsure{$T$|$\lambda$} \illegible{2} \unsure{$(20)$}} par
\[
\deleted{\unsure{\lambda|T_m|T_n}}\ T_\lambda(z_1,z_2,z)=\left(\frac{1}{\lambda}z_1,\frac{1}{\lambda}z_2,\frac{1}{\lambda}z\right).\tag*{(22)}
\]
% slanted
\margin{opération de $\mathbb C^*$ définie par le pb modulaire, la base $\varphi$ \inserted{(de $t_E$)} multipliée par $\lambda\varphi$ \ldots}

Chacun des deux groupes \unsure{commutants|commutant|commute} $\mathbb C^*,SL(2,\mathbb Z)$ opère librement sur $X_{PD^+}$ et même sur $PD^+$, mais l'opération du groupe produit $\mathbb C^*\times SL(2,\mathbb Z)$ \inserted{sur $PD^+$} n'est pas une opération libre, p.ex. l'élément $(-1,-1)$ opère trivialement sur $PD^+$ (autom. universel des courbes elliptiques) -- mais non sur $X_{PD^+}$, mais il \deleted{\unsure{opère|n'a pas}} a des pts fixes \deleted{\unsure{à savoir|de moi}} $X_{PD^+}$ (et d'autres ?!)

\deleted{En fait, l'opération sur $X_{PD^+}$ est une opération libre, sauf en les \unsure{T|I}}

\deleted{quotient par $\mathbb C^*$ étant $X_{D^+}$, et le quotient de ce dernier par $SL(2,\mathbb Z)$, \inserted{moins la section nulle} s'identifiant à $M_{1,2}$ !}
% archive sheet number: 151
\scanpage{153}
\begin{flushright}
normalité des ext. quadratiques
\end{flushright}
% archive sheet number: 152
\scanpage{154}
\deleted{
\diagram{\begin{tikzcd}[ampersand replacement=\&,row sep=small,column sep=small]
\& \& \unsure{G_0|G_\infty|\mathfrak S_0} \& \& \\
1\arrow[r] \& \deleted{\unsure{A|H}}N\arrow[r] \& G\arrow[r]\arrow[u,hook] \& \unsure{G_A|\mathfrak S_A}\arrow[r,dash]\arrow[ul,bend right] \& 1\\
\& ((\mathbb F_2)^A)' \arrow[u,equal] \& \& \unsure{G_A|\mathfrak S_A}\arrow[u,equal] \& \\
\& \ker(\mathbb F_2^A\mathbin{\unsure{-|\to}}\mathbb F_2)\arrow[u,equal] \& \& \&
\end{tikzcd}}
\[
(x_i)\mathrel{\unsure{\mapsto|\vdash}}\sum x_i
\]
\[
\begin{aligned}
ux&=yu & y&=uxu^{-1}\\
&&&=u(x)
\end{aligned}
\]
\[
\begin{aligned}
g&=\underset{\substack{\uparrow\\G_0}}{u}\cdot\underset{\substack{\uparrow\\N}}{x}&&=yu\\
g'&=u'\cdot x'&&=y'u'
\end{aligned}
\qquad
g'g=y'\underbrace{u'yu'^{-1}}\underbrace{u'u}
\]
\[
\begin{aligned}
g'g&=u'x'ux=(u'u)(\underbrace{u^{-1}x'u}_{u^{-1}(x')x})x
=\underset{\substack{\Vert\\u'u}}{\unsure{U|u}''}\underset{\substack{\Vert\\u^{-1}(x')x}}{x''}\\
&=(y'\deleted{\unsure{y}}u'(y))u'u
\end{aligned}
\]
}
\medskip\hrule\medskip

L'extension quadratique de $\mathbb Q(t)$ est induite par l'extension quadratique
\[
B=\mathbb Z[t]/(S^2-P(t))
\]
de $A=\mathbb Z[T]$. On voit que pour $\unsure{\forall|\mathrm{tt}}\,p\in\mathbb P\setminus\{2\}$, localisant $A$ en $pA$ d'où $A_{pA}=\mathbb Z_{p\mathbb Z}\{t\}\overset{\mathrm{dfn}}=A(p)$, $\deleted{\unsure{B_b|B_6}}B(p)\overset{\mathrm{dfn}}=B\otimes_A A(p)$ est normal (que $p$ divise $q$, ou non). Si $p=2$, alors $B(2)$ est normal si $\deleted{\unsure{p}}2\mid q$. Si $\deleted{\unsure{p}}2\nmid q$, alors il est normal aussi si $Q\bmod 2\in\mathbb F_2[t]$ n'est pas un carré dans $\mathbb F_2[t]$ (donc pas non plus dans $\mathbb F_2(t]$). \deleted{\unsure{S'il|Si}} Il en est de $\overline{\overline{\phantom{m}}}$ s'il est un carré dans $\mathbb F_2[t]$, mais pas dans $\mathbb Z/4\mathbb Z[t]$ -- i.e un certain invariant dans $\mathbb Z/2\mathbb Z[t]$ n'est pas nul. Donc le seul cas où $B(2)$ ne soit normal est celui où $\deleted{\unsure{p}}2\nmid q$ \deleted{\unsure{et|$Q$}} $Q(t)$ est un carré mod $4$ (p.ex. $Q(t)=t^2-4=(T+2)(T-2)$ \inserted{ou $Q(t)=4t+1,\ \varepsilon=+1$ ou $Q(t)\unsure{+|=}4t-1,\ \varepsilon=-1$}).

Si on excepte ce cas, on voit que $B$ est normal en codim $1$'' -- i.e \inserted{$\mathrm{Spec}\,B$} n'a qu'un nb fini \deleted{d'idéaux premiers}\inserted{de pts non} normaux, et ces pts sont fermés (donc isolés). Donc \deleted{\unsure{en}} sauf \deleted{\unsure{l'est}} au dessus d'un sous-schéma fini de $\mathrm{Spec}\,A=(\mathrm{Spec}\,\mathbb Z)\unsure{[T]|[t]|/}$, $B$ représente \unsure{le|la} normalisé de $\mathbb E^1_{\mathbb Z}$ dans l'extension quadratique envisagée -- et ceci vaut encore le cas sans restriction, si on remplace $\mathbb Z$ par $\mathbb Z[1/2]$.
% archive sheet number: 153
\scanpage{155}
Normalité des extensions quadratiques \quad $A[S]/(S^2-a)=B$

On suppose d'abord $A$ anneau jauge, soit $p$ sa car. résiduelle.
% slanted
\margin{\uline{NB} Si $p\ne2$, alors $B$ est local sss \uline{ou} $=a\in\mathfrak m$ \uline{ou} ($a\deleted{\in}\notin\mathfrak m$ et) \deleted{\unsure{$a\in$}} $\dot a\notin k^2$}
$p\ne2$, alors $B$ normal sss $a\notin\mathfrak m^2$ -- et cette condition est nécessaire en tous cas (y compris $p=2$). Si $p=2$, distinguons deux cas:

a) $a\in\mathfrak m$ alors $B$ normal sss $a\notin\mathfrak m^2$

b) $a\notin\mathfrak m$ i.e. $u$ une unité. Soit $\dot a$ son image dans $K'=k(A)=A/\mathfrak m$.

$1^\circ$) $\dot a\notin k^2$, alors $B$ est normal (c'est un anneau jauge de corps résiduel $k[\sqrt{\dot a}]$)

$2^\circ$) $\dot a\in k^2$, soit $\dot a=\dot\alpha^2$, où $\alpha\in A$ déterminé mod $\deleted{\unsure{\mathfrak m}}\unsure{\mathfrak m}$ -- donc $a=\alpha^2\deleted{\unsure{+}}-b$, avec $b\inserted{=\alpha^2-a}\in\mathfrak m$. Posons $T=S-\alpha$, alors
% slanted
\margin{\uline{NB} $\widetilde b\deleted{\unsure{\in}}\in\mathfrak m/\mathfrak m^2$ ne dépend que de $a$, et pas du choix de $\alpha$.}
\[
\begin{aligned}
S^2-a&=(T+\alpha)^2-a=T^2+2\alpha T+b\\
B&\simeq A[T]/(T^2+2\alpha T+b),\quad
\widehat B\simeq\underbrace{A[[T]]}_{C}/(T^2+2\alpha T+b)
\end{aligned}
\]
% slanted
\margin{NB $2\in\mathfrak m\subset\mathfrak r(C)$ \quad $T\in\mathfrak r(C)$}
$T^2,2\alpha T$ sont dans $\deleted{\unsure{\mathfrak m|M}}\mathfrak r(C)^2$, donc $T^2+2\alpha T+b\notin(\mathfrak r(C))^2$ sss $b\notin\mathfrak r(C)^2$ i.e $b\notin\mathfrak m^2$, i.e. $\widetilde b\ne0$ dans $\mathfrak m/\mathfrak m^2$

Donc les cas de normalité sont les suivants
\[
\begin{aligned}
&a\in\mathfrak m\setminus\mathfrak m^2,\ \text{\uline{ou}}\
a\ \text{de la forme }\alpha^2-b\ \text{avec }\alpha\ \text{unité},\
b\in\mathfrak m\setminus\mathfrak m^2\\
&\text{\uline{ou}}\quad a\notin\mathfrak m,\quad\dot a\notin k^2
\end{aligned}
\]
% slanted
\margin{Donc les cas où $B$ \uline{non} normal sont ceux où $a\in\mathfrak m^2$, \uline{ou} $\unsure{\widehat a|\dot a|a}\bmod\mathfrak m^2$ le carré d'une unité de $A/\mathfrak m^2$. \uline{Ce dernier cas \unsure{correspond|correspondra} à une singularité quadratique ordinaire.}}

\uline{Application} aux rev. quadratiques de $\mathbb P^1_{\mathbb Z}$ \inserted{(ou de $\mathbb E^1_{\mathbb Z}$)} génériquement \deleted{\unsure{non|normales}} étales sur chaque fibre \inserted{au} dessus de $\mathrm{Spec}\,\mathbb Z$.

Le revêtement \deleted{\unsure{au dessus du pt}}\inserted{quadratique} du pt gén. correspond à une ext. quadratique de $\mathbb Q(t)$, donc décrit par un élément de $\mathbb Q(t)^*/\mathbb Q(t)^{\unsure{*|2|*2}}$.

Or, $\mathbb Q(t)^{*}/\unsure{\{\pm1\}}$ a une base formée des éléments suivants

a) Les nbs premiers $p\in\mathbb P$ \unsure{!} $(\mathbb P=\{2,3,5,7\ldots\})$

b) Les polynomes $P=a_0t^n+a_1t^{n-1}+\ldots+a_n\in\mathbb Z[t]$, avec: $a_0>\deleted{\unsure{1}}0$, \deleted{\unsure{et}} les $a_i$ premiers entre eux (i.e. $\deleted{\unsure{\forall}}\forall p\in\mathbb P$, $P\bmod p\in\mathbb Z_p[t]$ \deleted{est} $\ne0$), $P$ premier dans $\mathbb Q[T]$.

Ainsi un élément de $\mathbb Q\unsure{[t]|(t)|(t)^*}/\mathbb Q(t)^{\unsure{*|2|*2}}$ est représenté de façon unique par un produit
\[
\unsure{P|\mathbb P}=\underbrace{\pm}_{\varepsilon}\,
\underbrace{p_1p_2\cdots p_{\unsure{h|k|n}}}_{q}\,
\underbrace{P_1\cdots P_\ell}_{Q}\tag*{(1)}
\]
où les $p_i\in\mathbb P$ sont distincts, et les $P_j$ comme dans b) aussi ! Cela signifie ainsi qu'on aura
\[
P=\pm qQ\qquad
\left\{\begin{aligned}
&q\in\mathbb N\ \text{``quadratfrei''}\\
&Q\in\mathbb Z[t]\ \text{à coeff. dominant }>0\\
&\text{``quadratfrei'' dans }\mathbb Q[t]\\
&\text{à coeff. premiers entre eux.}
\end{aligned}\right.
\]
% archive sheet number: 154
\scanpage{156}
Considérons \unsure{d'autre|d'entrée|d'abord} \unsure{par|part|pour} $B\otimes\mathbb F_p$ sur $A\otimes\mathbb F_p=\mathbb F_p[T]$, donné par l'équation \quad $S^2-P(T)=0$ \deleted{\unsure{en car.}} sur $\mathbb F_p$. \inserted{Supposons $p\ne2$,} Cette extension est \uline{gén.}\textsuperscript{t} \deleted{\unsure{non ramifiée}}\inserted{étale} \deleted{\unsure{pour}}\inserted{sss} $P(T)\not\equiv0\ p$. Cela montre déjà que \deleted{\unsure{des} \unsure{$q$}} l'on doit avoir dans (1) $q\in\{1,2\}$, et dans le cas où $B(\deleted{\unsure{p}}p)$ est aussi normal \inserted{(p.ex. si $q=2$)}, et où on veut donc que $B\otimes\mathbb F_2$ soit gén. étale sur $A\otimes\mathbb F_2$, on voit que l'on doit avoir $q=1$. Mais l'équation $S^2-P(T)\bigl(=S^2-\varepsilon Q(T)\bigr)=0$ en car $2$ n'est jamais gén. étale ! Donc \deleted{il faut} pour avoir qque chose gén. étale en \uline{$\unsure{\forall|\mathrm{tt}}$ cas}, il \uline{faut} que $B(p)$ soit non normal i.e. \inserted{que} $P=\underset{\substack{\uparrow\\\{\pm1\}}}{\varepsilon}\deleted{\underset{\{1,2\}}{\unsure{q}}}Q$ soit un carré mod $4$. \uline{NB} Pour $Q$ fixé, ceci fixe le signe $\varepsilon$ sans ambiguité, car $-1$ n'est pas un carré mod $4$ !). Cette condition est aussi suffisante \ldots.

\uline{Exemple} (\deleted{\unsure{étant}}lié \unsure{aux} modules des courbes elliptiques). Soit $Q(T)=T(T-\lambda)$, avec $\lambda\in\mathbb Z\setminus\{0\}$ (dans le cas qui m'intéresse, \deleted{$\lambda$} $\lambda=2^\alpha3^\beta$ l'invariant modulaire exceptionnel). On a $Q(T)\equiv T^{\unsure{2|3}}\deleted{\unsure{\equiv0}}(2)$ \unsure{si|ssi} $\lambda\equiv0\ (4)$. On trouve un \unsure{revêt\textsuperscript{t}|revêt} étale de
\[
\mathbb E^1_{\mathbb Z}\setminus\{0,\lambda\}=\mathrm{Spec}\,\mathbb Z[T]_{T(T-\lambda)}
\]
et $\overline{\overline{\phantom{m}}}$ de
\[
\mathbb P^1_{\mathbb Z}\setminus\{0,\lambda\}=\mathrm{Spec}\,\mathbb Z[T']_{\lambda T'-1}
\]
% slanted
\margin{on a fait $T'=1/T$}
en normalisant $\mathbb Z[T']_{(\lambda T'-1)}$ dans $\mathbb Z[T']_{(\lambda T'-1)}/(\unsure{\delta|s|S}^2-(\lambda T'-1))$
% archive sheet number: 155
\scanpage{157}
% letterhead
\begin{tabular}{@{}c@{\qquad}c@{}}
UNIVERSITÉ DES SCIENCES &\\
ET TECHNIQUES DU LANGUEDOC & ACAD\unsure{É|E}MIE DE MONTPELLIER\\
MONTPELLIER II &\\
Place Eugène Bataillon &\\
34060 MONTPELLIER CEDEX &
\end{tabular}

Tel. : 63.91.44 -- Télex : USTMONT 49944 F

Service de la Scolarité\\
MD/gs

Montpellier, le 9 JUIN 1981 \hfill 19\\
\[
\handwritten{T^{\unsure{2|\prime}}+2\alpha T+b}
\]

M \,M \quad Le Professeur \quad \handwritten{GROTHENDIECK}\\
\[
\handwritten{T^2+2\alpha T+b}
\]
% printed form
\uline{OBJET} : \emph{Surveillance des examens}
\[
\handwritten{T^2+}
\]
\[
\handwritten{\mathbb F_2[T,\pi]\qquad T^2+\alpha\pi T+\beta\pi^2\,,}
\]
\[
\begin{aligned}
\handwritten{x^2+\alpha xy\unsure{+|\alpha}\beta^2 y^2}\ \text{\scriptsize M}
&\qquad \handwritten{\alpha^2-4\beta=\alpha^2}\\
\handwritten{(x+\deleted{\unsure{\alpha|b|\surd}}\beta y}
&\qquad \handwritten{\unsure{-1|\cdot\,1|1}}
\end{aligned}
\]

J'ai l'honneur de rappeler à votre attention que vous avez été désigné(e) pour assurer la surveillance des épreuves des examens qui doivent avoir lieu à l'U.S.T.L.
\[
\handwritten{4T\unsure{-|=}1}
\]

le \qquad : \quad 16 JUIN 1981\\
de \qquad : \quad \handwritten{14} h. \handwritten{45} à \handwritten{18} h.\\
Amphi \quad : \quad \handwritten{9 et 10}

\noindent\fbox{\parbox{.91\linewidth}{\textbullet\ Si pour diverses raisons il ne vous était pas possible d'assurer cette surveillance, il vous appartiendrait de vous trouver un remplaçant.}}
% archive sheet number: 156
